\subsection{KT-R2-POINT and the actual geometric gate}
\label{sec:kt-r2-point}

\begin{lemma}[KT-R2-POINT]
For the literal question in item~7, over
$\kappa=\overline{\mathbf F}_2$, the chosen point $P$ imposes
no extra condition on existence. A connected finite \'etale
cyclic cover $D\to E$ with primitive factor of two-rank zero
can be pointed above every $P\in E(\kappa)$.
\end{lemma}

\begin{proof}
A finite connected cover is surjective. Its fibre over $P$ is
a nonempty finite \'etale $\kappa$-scheme, hence a disjoint
union of copies of $\operatorname{Spec}\kappa$. Choosing one
of these points gives the pointing. Neither the cover nor its
primitive abelian factor is changed. All equations, maps,
$P$ and the chosen lift descend to some common finite field;
geometric two-rank is unchanged by this descent or by a
further extension of constants.
\end{proof}

This removes the parameter $P$ from the \emph{geometric}
existence question. It does not remove the prescribed-field
pointing condition in AN\_GRH A45--A52: those results measure
the rational two-primary group over a fixed finite field.

\paragraph{An exact genus-two theta-orbit formulation.}
Let $E/\kappa$ have genus two, let $L\in\operatorname{Pic}^0(E)$
have exact odd order $n>1$, and let $D\to E$ be its cyclic
cover. Write
\[
 0\longrightarrow\mathcal O_E\longrightarrow
 F_*\mathcal O_E\longrightarrow B_E\longrightarrow0
\]
for the locally exact differential sequence, using absolute
Frobenius and its semilinear action. Since the characteristic
is two, $B_E$ is a line bundle of degree one. Put
\[
 \Theta_{B_E}=\{M\in\operatorname{Pic}^0(E):
                         H^0(E,B_E\otimes M)\ne0\}.
                                                               \tag{R2.1}
\]
This is a translate of the Abel--Jacobi image of $E$.
The remaining existence question is exactly to find such
an odd-order $L$ for which every orbit of multiplication by
two on $(\mathbf Z/n\mathbf Z)^\times$ contains an $i$
with $L^i\in\Theta_{B_E}$.

Here is a proof of this particular form of the standard
Raynaud criterion. The character decomposition of
$H^1(D,\mathcal O_D)$ has summands $H^1(E,L^i)$.
For $i\ne0\pmod n$, Riemann--Roch and $H^0(E,L^i)=0$
give $\dim H^1(E,L^i)=1$. Tensoring the locally exact
differential sequence by $L^i$, then applying the projection
formula and coherent cohomology, identifies
\[
 \ker\bigl(F:H^1(E,L^i)\to H^1(E,L^{2i})\bigr)
       =H^0(E,B_E\otimes L^i).                              \tag{R2.2}
\]
The map is semilinear; each arrow between these one-dimensional
spaces is either zero or a semilinear isomorphism.
On a complete doubling orbit, stable Frobenius rank is zero
precisely when at least one arrow is zero. The primitive
factor selects the characters with $\gcd(i,n)=1$; because
$n$ is odd, its idempotent is integral at two and computes
this same stable part. Therefore, if
$h_n=\operatorname{ord}_n(2)$,
\[
 f_2(P_n)=h_n\cdot
 \#\{\text{doubling orbits of primitive }i:
                   L^i\notin\Theta_{B_E}\text{ for all }i
                   \text{ in the orbit}\}.                 \tag{R2.3}
\]
The identification of stable coherent Frobenius rank with
two-rank and the character version of the locally exact
differential criterion are the standard inputs; compare
Pop--Sa\"idi (2003), \S1.3, p.~112, carried in AN\_GRH A48,
and Sa\"idi, \emph{On complete families of curves with a
given fundamental group in positive characteristic},
Manuscripta Math.~118 (2005), \S\S1--2, pp.~428--431.
The displayed formulation is an explicit proof of the form
needed here, not a new general Raynaud theorem.

\paragraph{What was and was not advanced.}
The Rosati-fixed $\mathcal O_5$ multiplication remains an
extra restriction on the allowed genus-two curves. It is not
used in (R2.1)--(R2.3), and no argument here forces the required
orbit incidences on that restricted locus. The known failures
AN\_GRH A36--A44 and the sharper finite-field alternatives
A45--A52 are retained as supplied background. Searches of the
Raynaud--Sa\"idi--Tamagawa and Ozman--Pries literature did not
supply the missing all-curves incidence theorem. In particular,
Ozman--Pries, \emph{Ordinary and almost ordinary Prym varieties},
Theorem~1.1, concerns ordinary Pryms over a generic curve in a
fixed $p$-rank stratum; it is not an existence theorem for
two-rank-zero primitive factors over every prescribed
$\overline{\mathbf F}_2$-curve. No universal exponent or
annihilator-index bound in the actual primitive rational group
has been obtained, and no curve for which every odd cyclic
cover fails has been exhibited.

